\section*{Bab \ref{ch:b1:sinusoidal-impedance} --- Keadaan Tunak Sinusoidal dan Impedansi}

\begin{solution}{exo:b1:sinusoidal-impedance:1}
$\underline U = 5\eu^{j\pi/3} = 2.5 + 4.33j$ (V). $3\cos\omega t +
4\sin\omega t$: $3 + 4\eu^{-j\pi/2} = 3 - 4j$, modulusnya $5$, argumennya
$-\arctan(4/3) = \qty{-0.927}{rad}$: $5\cos(\omega t - 0.927)$.
\end{solution}

\begin{solution}{exo:b1:sinusoidal-impedance:2}
$R$: \qty{100}{\ohm} pada keduanya. $L\omega$: \qty{31.4}{\ohm} (\qty{50}{Hz}),
\qty{628}{\ohm} (\qty{1}{kHz}). $1/C\omega$: \qty{318}{\ohm}, \qty{15.9}{\ohm}.
$RL$ seri pada \qty{50}{Hz}: $\abs{\underline Z} = \sqrt{100^2 + 31.4^2} =
\qty{105}{\ohm}$, fasenya $\arctan(0.314) = 17^\circ$.
\end{solution}

\begin{solution}{exo:b1:sinusoidal-impedance:3}
$U = 230\sqrt2 = \qty{325}{V}$. $I_{\mathrm{rms}} = 2000/230 = \qty{8.7}{A}$,
puncaknya \qty{12.3}{A}. $\langle Ri^2\rangle = R\langle i^2\rangle =
RI_{\mathrm{rms}}^2$ menurut definisi nilai rmsnya.
\end{solution}

\begin{solution}{exo:b1:sinusoidal-impedance:4}
$f_0 = 1/(2\pi\sqrt{LC}) = \qty{5.03}{kHz}$; $Q = \sqrt{L/C}/R = 3.16$;
$\Delta f = f_0/Q = \qty{1.6}{kHz}$; $I_{\max} = E/R = \qty{10}{mA}$.
\end{solution}

\begin{solution}{exo:b1:sinusoidal-impedance:5}
$1/\sqrt2$ pada $\omega = 1/RC$: $f_c = 1/(2\pi RC) = \qty{1.0}{kHz}$. Pada
\qty{1.0}{kHz}: nisbahnya $0.71$, fasenya $-45^\circ$. Pada \qty{10}{kHz}:
$RC\omega = 10$, nisbahnya $1/\sqrt{101} = 0.10$, fasenya $-84^\circ$.
\end{solution}

\begin{solution}{exo:b1:sinusoidal-impedance:6}
$\underline Z = 100 + j(250 - 100) = 100 + 150j$: $Z = \qty{180}{\ohm}$,
$\arg = 56^\circ$. $I = 1.0/180 = \qty{5.6}{mA}$, tertinggal di belakang
sumbernya sebesar $56^\circ$. $U_R = \qty{0.56}{V}$, $U_L = \qty{1.39}{V}$, $U_C = \qty{0.56}{V}$;
\omterm{def:b1:sinusoidal-impedance:complex}{fasornya}: $U_R$ searah $I$, $U_L$ ke atas, $U_C$ ke bawah, resultannya \qty{1.0}{V}.
$L\omega > 1/C\omega$: di atas resonansi (induktif).
\end{solution}

\begin{solution}{exo:b1:sinusoidal-impedance:7}
$f_0 = 1/(2\pi\sqrt{LC})$: \qty{1.59}{MHz} (\qty{50}{pF}) sampai \qty{0.50}{MHz}
(\qty{500}{pF}). Pada \qty{1.00}{MHz}: $C = 1/(4\pi^2f_0^2L) = \qty{127}{pF}$.
$\Delta f = \qty{10}{kHz}$. Stasiun pada \qty{1.009}{MHz}: $x = 1.009$,
$Q(x - 1/x) = 1.8$, nisbahnya $1/\sqrt{1 + 3.2} = 0.49$. \omterm{def:b1:dc-circuits:voltage}{Tegangan}
kapasitornya $QE = \qty{1.0}{mV}$.
\end{solution}

\begin{solution}{exo:b1:sinusoidal-impedance:8}
$Q(x - 1/x) = \pm1 \Leftrightarrow x^2 \mp x/Q - 1 = 0$; akar positifnya
$x_{1,2} = \mp 1/2Q + \sqrt{1 + 1/4Q^2}$. Selisihnya $1/Q$, sehingga
$\Delta\omega = \omega_0/Q$; hasil kalinya $(1 + 1/4Q^2) - 1/4Q^2 = 1$, sehingga
$\omega_1\omega_2 = \omega_0^2$ (pitanya setangkup pada skala logaritmik).
\end{solution}

\begin{solution}{exo:b1:sinusoidal-impedance:9}
$I = P/(U\cos\varphi) = 2000/(230 \times 0.70) = \qty{12.4}{A}$. \omterm{thm:b1:sinusoidal-impedance:power}{Daya
reaktifnya} $P\tan\varphi = 2000 \times 1.02 = \qty{2.0}{kvar}$. $C = 2040/(314
\times 230^2) = \qty{123}{\micro F}$. Arus barunya $2000/230 = \qty{8.7}{A}$.
Ruginya $0.50 \times 12.4^2 = \qty{77}{W}$ sebelumnya, $0.50 \times 8.7^2 =
\qty{38}{W}$ sesudahnya.
\end{solution}

\begin{solution}{exo:b1:sinusoidal-impedance:10}
$\underline Y = 1/R + j(C\omega - 1/L\omega)$; $\underline U = \underline I/
\underline Y$ maksimum ketika suseptansinya lenyap, yaitu pada $\omega_0$,
tempat $U = RI$. Setengah daya: $\abs{C\omega - 1/L\omega} = 1/R$, yaitu
$RC\omega_0(x - 1/x) = \pm1$: $\Delta\omega = \omega_0/(RC\omega_0) = 1/RC$
dan $Q = RC\omega_0 = R\sqrt{C/L}$. Angkanya: $\omega_0 = \qty{1.0e6}{rad/s}$
($f_0 = \qty{159}{kHz}$), $Q = 10^4\sqrt{10^{-6}} = 10$, $\Delta f =
\qty{16}{kHz}$.
\end{solution}

\begin{solution}{exo:b1:sinusoidal-impedance:11}
$U_C = I/C\omega = E/[C\omega\sqrt{R^2 + (L\omega - 1/C\omega)^2}]$;
isinya dikalikan $C\omega$: $\sqrt{(RC\omega)^2 + (LC\omega^2 - 1)^2} =
\sqrt{x^2/Q^2 + (1 - x^2)^2}$. Dengan $s = x^2$: $(1 - s)^2 + s/Q^2$
minimum pada $s = 1 - 1/2Q^2$ (bila positif), minimumnya $1/Q^2 - 1/4Q^4$:
$U_{C,\max} = QE/\sqrt{1 - 1/4Q^2}$. $Q = 3.16$: $x_r = 0.975$,
$U_{C,\max} = 3.2E$. Untuk $Q < 1/\sqrt2$ minimumnya berada di $s = 0$: $U_C$
menurun secara monoton dari $E$, tanpa puncak.
\end{solution}

\begin{solution}{exo:b1:sinusoidal-impedance:12}
Pada resonansi $i = I\cos\omega_0 t$ dan $u_C = (I/C\omega_0)\sin\omega_0 t$,
sehingga, dengan memakai $1/C\omega_0^2 = L$,
\[
\tfrac12 Li^2 + \tfrac12 Cu_C^2 = \tfrac12 LI^2\cos^2\omega_0 t
+ \tfrac12 LI^2\sin^2\omega_0 t = \tfrac12 LI^2 .
\]
Dilesapkan tiap periode: $\tfrac12 RI^2T$. Nisbahnya $\times 2\pi$: $2\pi L/(RT)
= L\omega_0/R = Q$.
\end{solution}

\begin{solution}{pb:b1:sinusoidal-impedance:1}
\textbf{1.} $\underline Z = R + j(L\omega - 1/C\omega)$.

\textbf{2.} $\underline I = E/\underline Z$; $I = E/\sqrt{R^2 + (L\omega -
1/C\omega)^2}$; $\varphi_i = -\arctan[(L\omega - 1/C\omega)/R]$.

\textbf{3.} \omterm{def:b1:sinusoidal-impedance:impedance}{Reaktansinya} lenyap pada $\omega_0 = 1/\sqrt{LC}$: $I_{\max}
= E/R$, $\varphi_i = 0$.

\textbf{4.} $C = 1/(4\pi^2f_0^2L) = \qty{127}{pF}$; $L\omega_0 = 2\times10^{-4}
\times 6.28\times10^6 = \qty{1257}{\ohm}$.

\textbf{5.} $Q = 1257/12.6 = 100$.

\textbf{6.} $L\omega - 1/C\omega = L\omega_0(x - 1/x)$, bagi dengan $R$.

\textbf{7.} $x < 1$: \omterm{def:b1:sinusoidal-impedance:impedance}{reaktansinya} negatif, $\varphi_i > 0$, arusnya
mendahului (kapasitif); $x > 1$: tertinggal (induktif).

\textbf{8.} $Q(x - 1/x) = \pm1$.

\textbf{9.} $x^2 \mp x/Q - 1 = 0$, $x_{1,2} = \mp1/2Q + \sqrt{1 + 1/4Q^2}$,
$x_2 - x_1 = 1/Q$: $\Delta\omega = \omega_0/Q$.

\textbf{10.} $\Delta f = \qty{10}{kHz}$.

\textbf{11.} \qty{1.010}{MHz}: $Q(x - 1/x) = 100 \times 0.0199 = 2.0$,
nisbahnya $1/\sqrt{5.0} = 0.45$, dayanya $0.20$: \qty{-7}{dB}. \qty{1.020}{MHz}:
$3.96$, nisbahnya $0.25$, dayanya $0.06$: \qty{-12}{dB}.

\textbf{12.} Pas-pasan: tetangganya hanya \qty{7}{dB} lebih rendah. $Q$
yang lebih tinggi ($R$ yang lebih rendah), atau beberapa tingkat
terlaras yang berderet, mempertajam pemilihannya.

\textbf{13.} \qty{0.50}{MHz} sampai \qty{1.59}{MHz}: pita modulasi \omterm{def:b1:wave-propagation:signal}{amplitudo}.

\textbf{14.} $\underline U_C = \underline I/jC\omega$; pada $\omega_0$,
$U_C = I_{\max}/C\omega_0 = E/(RC\omega_0) = QE$.

\textbf{15.} $I_{\max} = 10^{-5}/12.6 = \qty{0.79}{\micro A}$; $U_C = QE =
\qty{1.0}{mV}$: penguatan \omterm{def:b1:dc-circuits:voltage}{tegangan} seratus kali tanpa penguat apa pun.

\textbf{16.} $\underline U_L = jL\omega_0\underline I = jQE$, $\underline U_C
= -jQE$: sama besar dan berlawanan, keduanya meniadakan; sumbernya melihat $R$ saja.

\textbf{17.} Tersimpan $\tfrac12 LI_{\max}^2 = \qty{6.2e-17}{J}$; dilesapkan
tiap periode $\tfrac12 RI_{\max}^2T = \qty{3.9e-18}{J}$; $2\pi \times 6.2/0.39
\approx 100 = Q$.

\textbf{18.} $\tau = 2Q/\omega_0 = \qty{32}{\micro s}$. Rangkaiannya tak
dapat mengikuti perubahan yang lebih cepat daripada $\sim\tau$: \omterm{thm:b1:sinusoidal-impedance:resonance}{lebar
pita} $\Delta f$ hanya meloloskan modulasi yang lebih lambat daripada
$\Delta f$; dengan $Q = 1000$, $\Delta f =
\qty{1}{kHz}$ dan nada tinggi acaranya akan hilang.

\textbf{19.} $Q$ berkurang separuh menjadi $50$, $\Delta f$ berlipat dua
menjadi \qty{20}{kHz}, $U_C$ berkurang separuh menjadi \qty{0.5}{mV}.

\textbf{20.} $R_{\text{tot}} = \qty{62.6}{\ohm}$: $Q = 1257/62.6 = 20$,
$\Delta f = \qty{50}{kHz}$, $U_C = 20 \times \qty{10}{\micro V} = \qty{0.2}{mV}$:
kepekaan dan selektivitasnya sama-sama runtuh.

\textbf{21.} $P = E^2/(2R_{\text{tot}}) = 10^{-10}/125 = \qty{8e-13}{W}$;
bagian di $R_a$: $50/62.6 = 80\%$.

\textbf{22.} Penyesuaian menuntut hambatan beban yang sama dengan $R_a$,
yaitu hambatan total yang besar; selektivitas menuntut hambatan total
sekecil mungkin ($Q = L\omega_0/R_{\text{tot}}$). Orang tak dapat
memperoleh keduanya dengan sambungan langsung.

\textbf{23.} Sebuah sadapan atau kapasitor gandeng yang kecil hanya
menyodorkan sebagian rangkaian resonannya kepada antena: hambatan
antenanya ``terlihat'' jauh berkurang oleh rangkaiannya, sehingga $Q$
tetap tinggi, sementara antenanya tetap memberi daya di dekat optimumnya.

\textbf{24.} Pada frekuensi setengah daya \omterm{def:b1:sinusoidal-impedance:impedance}{reaktansinya} sama dengan
$\pm R$, sehingga $\varphi = \pm45^\circ$ dan $\cos\varphi = 1/\sqrt2$;
dengan $I =
I_{\max}/\sqrt2$ pula, $P = \tfrac12 EI\cos\varphi$ menjadi separuh
nilainya pada resonansi.

\textbf{25.} $Q$ mengalikan \omterm{def:b1:wave-propagation:signal}{sinyalnya} dengan $Q$ dan menyempitkan pitanya
menjadi $f_0/Q$, dengan harga tanggapan yang lebih lambat
($\tau = 2Q/\omega_0$) dan gandengan yang rapuh ke antenanya; untuk kanal
\qty{10}{kHz} pada \qty{1}{MHz}, $Q \approx 100$ adalah angkanya --- dan
antenanya harus digandengkan secara ringan agar angka itu bertahan.
\end{solution}
